\documentclass[11pt,reqno]{amsart}

\usepackage[margin=1in]{geometry}
\usepackage{amsmath,amssymb,amsthm,mathtools}
\usepackage{microtype}
\usepackage[T1]{fontenc}
\usepackage{lmodern}
\usepackage{enumitem}
\usepackage{hyperref}

\hypersetup{
colorlinks=true,
linkcolor=blue,
citecolor=blue,
urlcolor=blue
}

\theoremstyle{plain}
\newtheorem{theorem}{Theorem}[section]
\newtheorem{proposition}[theorem]{Proposition}
\newtheorem{lemma}[theorem]{Lemma}
\newtheorem{corollary}[theorem]{Corollary}

\theoremstyle{definition}
\newtheorem{definition}[theorem]{Definition}
\newtheorem{hypothesis}[theorem]{Hypothesis}

\theoremstyle{remark}
\newtheorem{remark}[theorem]{Remark}

\DeclareMathOperator{\supp}{supp}
\DeclareMathOperator{\cum}{cum}
\DeclareMathOperator{\Tr}{Tr}
\DeclareMathOperator{\spanop}{span}
\DeclareMathOperator{\LP}{\mathcal{LP}}
\DeclareMathOperator{\Sp}{Sp}
\DeclareMathOperator{\Mp}{Mp}
\DeclareMathOperator{\Ad}{Ad}
\newcommand{\Hc}{\mathcal H}
\newcommand{\Bc}{\mathcal B}
\newcommand{\Lc}{\mathcal L}
\newcommand{\Dc}{\mathcal D}
\newcommand{\R}{\mathbb R}
\newcommand{\C}{\mathbb C}
\newcommand{\E}{\mathbb E}
\newcommand{\hh}{\mathfrak h}
\newcommand{\spalg}{\mathfrak{sp}}
\newcommand{\slalg}{\mathfrak{sl}}

\title[Band-Limited Hardy \(Z\), Edgeworth Cumulants, and Laguerre Positivity]{%
Band-Limited Pullbacks of Hardy \(Z\),\\
Orthogonal Spectral Measures, Edgeworth Cumulants,\\
and Spectral Laguerre Positivity}

\author{Stephen Crowley}
\address{Dallas, Texas, USA}
\date{April 30, 2026}

\begin{document}

\begin{abstract}
The normalized pullback of Hardy's \(Z\)-function by the Riemann--Siegel
phase is a finite-power stationary vector with sine-kernel covariance
\(K(h)=\sin h/h\). Band-limitedness and the orthogonal spectral
representation \(X(u)=\int_{-1}^{1}e^{iu\xi}\Phi(d\xi)\) follow. The
stationary pullback extends to a real entire function of order one and
type one, and its generalized Laguerre functionals admit the spectral form
\[
L_n[X](u)=\tfrac12\int_{[-1,1]^2}(\xi-\eta)^{2n}e^{iu(\xi+\eta)}
\Phi(d\xi)\Phi(d\eta).
\]
Membership in the Laguerre--P\'olya class \(\LP\) is equivalent to
positivity of all \(L_n[X](u)\), \(n\ge1\), \(u\in\R\). The zero
correspondence between \(X\) and \(\zeta\) on the critical line is exact,
so RH is unconditionally equivalent to this spectral Laguerre positivity.
The Gaussian realization is quasi-free; the Laguerre kernels are symplectic
generators with metaplectic character
\(\det(I-2itA_{n,u})^{-1/2}\), giving the cumulant formula
\(\tfrac{(m-1)!}{2}\Tr(A_{n,u}^m)\).
\end{abstract}

\maketitle

\section{The normalized pullback and the covariance theorem}

\begin{definition}[Phase pullback]
Let
\[
\theta(t)=\arg\Gamma\!\left(\frac14+\frac{it}{2}\right)-\frac t2\log\pi
\]
be the Riemann--Siegel phase. On the real line where \(\Theta=\theta\) is
strictly increasing for \(t\) large, put
\[
u=\Theta(t),\qquad t=t(u)=\Theta^{-1}(u).
\]
The variance-normalized pulled-back Hardy signal is
\[
X(u)=\frac{Z(t(u))}{\sqrt{2\Theta'(t(u))}},
\qquad
Z(t)=e^{i\theta(t)}\zeta\!\left(\tfrac12+it\right).
\]
\end{definition}

\begin{theorem}[Sine-kernel covariance]\label{thm:sine}
\[
K(h):=\lim_{U\to\infty}\frac1U\int_0^U X(u+h)X(u)\,du=\frac{\sin h}{h},
\qquad K(0)=1.
\]
\end{theorem}

\begin{proof}
\emph{Step 1: Scaling.}
On \(t\in[T,2T]\),
\(\theta'(t)=\tfrac12\log(t/2\pi)+O(t^{-1})\); set
\(L:=\tfrac12\log(T/2\pi)\) and \(N:=\lfloor\sqrt{T/2\pi}\rfloor\), so
\(U=\Theta(2T)-\Theta(T)\asymp TL\). Under \(u=\Theta(t)\) with
\(t_h:=\Theta^{-1}(\Theta(t)+h)\), Taylor expansion gives
\(t_h=t+h/\Theta'(t)+O(h^2/(TL^2))\) and the prefactor is
\(\tfrac12(1+O(1/(tL)))\). Hence
\[
K(h)=\lim_{T\to\infty}\frac1{2TL}\int_T^{2T}Z(t_h)Z(t)\,dt+O(1/L).
\]

\emph{Step 2: Four signs.}
Write \(Z=S+\bar S\) with
\(S(t):=\sum_{n\le N}n^{-1/2}e^{i\theta(t)-it\log n}+R_1(t)\) and
\(\|R_1\|_{L^2[T,2T]}^2=O(T^{1/2}\log T)\) (Titchmarsh, Ch.~VII). The
product \(Z(t_h)Z(t)\) splits into four sign-pieces plus \(R_1\)-cross
terms; the latter contribute \(O(T^{3/4}(\log T)^{1/2})\), i.e.\
\(O(T^{-1/4}/L)\) after division by \(TL\).

\emph{Step 3: Mixed-sign via Montgomery--Vaughan.}
For \(S_h\bar S\) the phase becomes \(ih(1-\log n/L)+it\log(m/n)+O(1/L)\).
MV gives off-diagonal \(O(N)=O(T^{1/2})\), and the diagonal yields
\(TL\int_0^1 e^{ih(1-x)}\,dx+O(T)\); the conjugate mixed piece gives the
conjugate. The Riemann-sum rate is \(O(1/L)\) uniformly for \(h\) in
compact sets by Abel summation against \(\sum_{n\le y}1/n=\log y+\gamma+O(1/y)\).

\emph{Step 4: Same-sign via Atkinson.}
For \(S_hS\) the phase is \(2i\theta(t)+ih+O(1/L)\); the functional equation
\(\chi(\tfrac12+it)=e^{-2i\theta(t)}\) converts
\(e^{2i\theta(t)}\zeta^2\) into a \(\chi^{-1}\)-twisted second moment,
which by Atkinson's formula (Ivi\'c, Ch.~15) is \(O(T^{1/2}\log T)\),
i.e.\ \(O(T^{-1/2}\log T/L)\) after division by \(TL\).

\emph{Step 5.}
\(\frac1{2TL}\int_T^{2T}Z(t_h)Z(t)\,dt=\int_0^1\cos(h(1-x))\,dx+O(T^{-1/4}+L^{-1})\to\sin h/h\).
\end{proof}

\section{The cyclic Hilbert space and translations}

\begin{definition}
\(\Hc_X=\overline{\spanop\{X(\cdot+h):h\in\R\}}\) with
\(\langle f,g\rangle_W=\lim_U\tfrac1U\int_0^U f\bar g\,du\); cyclic vector
\(x_0=X\).
\end{definition}

\begin{proposition}[Translations are unitary]\label{prop:unitary}
\((U_af)(u)=f(u+a)\) is a strongly continuous unitary group on \(\Hc_X\),
and \(\langle U_hx_0,x_0\rangle_W=K(h)\).
\end{proposition}

\begin{proof}
Wiener-mean translation invariance gives the isometry; \(U_{-a}\) is its
inverse. Strong continuity is continuity of \(K\) plus cyclic density.
\end{proof}

\section{Band-limitedness}

\begin{theorem}[Spectral representation]\label{thm:spectral}
There exists a projection-valued measure \(E\) with \(U_h=\int e^{ih\xi}\,E(d\xi)\)
and \(\Phi(B):=E(B)x_0\) satisfies
\(\langle\Phi(A),\Phi(B)\rangle_W=\mu(A\cap B)\) with \(\mu(B)=\|E(B)x_0\|_W^2\);
moreover \(X(u)=\int e^{iu\xi}\Phi(d\xi)\).
\end{theorem}

\begin{theorem}[Band-limitedness]\label{thm:bandlimited}
\(\mu(d\xi)=\tfrac12\mathbf 1_{[-1,1]}(\xi)\,d\xi\), and
\(X(u)=\int_{-1}^{1}e^{iu\xi}\Phi(d\xi)\).
\end{theorem}

\begin{proof}
\(\sin h/h=\int_{-1}^{1}e^{ih\xi}d\xi/2\); Bochner uniqueness identifies
\(\mu\). For \(B\cap[-1,1]=\varnothing\), \(\|\Phi(B)\|_W^2=\mu(B)=0\).
\end{proof}

\section{Entire continuation, order and type}

This section builds the entire function on \(\C\) that extends the real-line
stationary process \(X\), records its order and type, and matches it to the
hypothesis class of the Csordas--Dimitrov generalized Laguerre criterion.

\begin{proposition}[Paley--Wiener extension]\label{prop:PW}
The \(\Hc_X\)-valued map
\[
X(z):=\int_{-1}^{1}e^{iz\xi}\,\Phi(d\xi),\qquad z\in\C,
\]
is entire in \(z\), and for every \(z\in\C\),
\[
\|X(z)\|_W\le e^{|\operatorname{Im}z|}\,\sqrt{\|\Phi([-1,1])\|_W^2}
=e^{|\operatorname{Im}z|}.
\]
In particular, the scalar function \(X\) on the real line extends to an
entire function of exponential type at most one, of order one and type
exactly one.
\end{proposition}

\begin{proof}
Entirety is differentiation under the spectral integral, legitimate by
\(\|\Phi\|_W<\infty\) on the compact support \([-1,1]\). The bound is
\[
\|X(z)\|_W^2
=\int_{-1}^{1}|e^{iz\xi}|^2\,\mu(d\xi)
=\int_{-1}^{1}e^{-2\xi\operatorname{Im}z}\tfrac{d\xi}{2}
\le e^{2|\operatorname{Im}z|}.
\]
Sharpness of type one: take \(z=iy\); then
\(\|X(iy)\|_W^2=\tfrac{\sinh 2y}{2y}\sim\tfrac{e^{2y}}{4y}\) as \(y\to+\infty\).
\end{proof}

\begin{proposition}[Realness on the real line]\label{prop:real}
The scalar function \(X\) is real on \(\R\) and therefore its entire
continuation satisfies \(X(\bar z)=\overline{X(z)}\).
\end{proposition}

\begin{proof}
\(X(u)\) is \(\R\)-valued because \(Z(t)\) is real for \(t\in\R\) (Hardy's
\(Z\)) and \(\Theta'>0\) makes the normalization real. Entire functions
agreeing with a real function on \(\R\) satisfy the Schwarz reflection
identity.
\end{proof}

\begin{corollary}[\(X\) is in the Hermite--Biehler setting]\label{cor:HB}
\(X\) is a real entire function of exponential type one in the
finite-order Laguerre--P\'olya ambient class. In particular the
Csordas--Dimitrov generalized Laguerre criterion applies to \(X\).
\end{corollary}

\begin{proof}
Order one and finite type place \(X\) in Levin's class of real entire
functions of exponential type, which is contained in the ambient
finite-genus class of Csordas--Dimitrov.
\end{proof}

\section{The Cram\'er inverse}

\begin{theorem}[Inverse for increments]\label{thm:inverse}
For \(B=(a,b]\),
\(\Phi(B)=L_W^2\text{-}\lim_R\tfrac1{2\pi}\int_{-R}^{R}\tfrac{e^{-iau}-e^{-ibu}}{iu}X(u)\,du\).
\end{theorem}

\begin{proof}
Truncated Fourier inverses \(P_R\mathbf 1_B\to\mathbf 1_B\) in
\(L^2(\mu)\); isometry gives the conclusion.
\end{proof}

\begin{proposition}[No point-frequency inverse]\label{prop:nopoint}
For \(\lambda\in(-1,1)\), \(I_T(\lambda)=\pi^{-1}\int_0^T e^{-iu\lambda}X\,du\)
satisfies \(\|I_T\|_W^2\sim T/\pi\) and is not Cauchy.
\end{proposition}

\section{Gaussian realization and quasi-freeness}

\begin{definition}
\(W:\Hc_X\to L^2(\Omega)\) isonormal: \(\E[W(f)W(g)]=\langle f,g\rangle_W\),
and \(\Phi_G(B):=W(\Phi(B))\).
\end{definition}

\begin{theorem}[Quasi-freeness]\label{thm:quasifree}
\(\cum(W(f_1),\ldots,W(f_r))=0\) for \(r\ge3\).
\end{theorem}

\begin{proof}
Gaussian CGF is quadratic; derivatives of order \(\ge3\) vanish.
\end{proof}

\section{Heisenberg/metaplectic structure on the band}

\subsection{Translation generator}

\(U_a=e^{iaP}\) with \(P\) unitarily equivalent to multiplication by
\(\xi\) on \(L^2([-1,1],\tfrac12 d\xi)\) under the isometry
\(\Phi\colon f\mapsto\int f\,d\Phi\). The algebra \(\spanop\{P\}\) is
abelian: \(e^{iaP}e^{ibP}=e^{i(a+b)P}\).

\subsection{Heisenberg--Weyl algebra}

\begin{definition}
On \(\Dc:=C^\infty_c((-1,1))\subset L^2([-1,1],\tfrac12 d\xi)\),
\(Q:=i\partial_\xi\).
\end{definition}

\begin{proposition}[CCR]\label{prop:ccr}
\([Q,P]=i\,\mathbf 1\) on \(\Dc\); the real span
\(\hh_3=\spanop_{\R}\{Q,P,\mathbf 1\}\) is the Heisenberg algebra.
\end{proposition}

\begin{theorem}[Weyl relation]\label{thm:weyl}
\(e^{i(aP+bQ)}=e^{-iab/2}e^{iaP}e^{ibQ}\); BCH on \(\hh_3\) terminates because
\([\hh_3,\hh_3]=\R\mathbf 1\) is central.
\end{theorem}

\begin{remark}
\(Q\) has deficiency \((1,1)\); any self-adjoint extension from
\(f(-1)\mapsto e^{i\varphi}f(1)\) gives a representation of the Heisenberg
group. Downstream statements depend only on the infinitesimal CCR.
\end{remark}

\subsection{Symplectic Lie algebra}

\begin{proposition}[Quadratics are symplectic]\label{prop:spalg}
For \(A=A^*\) on \(\Hc\), \(Q_A=\tfrac12\langle G,AG\rangle\) satisfies
\([Q_{A_1},Q_{A_2}]=Q_{i[A_1,A_2]}+(\text{central})\). The map
\(A\mapsto Q_A\) embeds into \(\spalg(\Hc^\R)\), and \(e^{itQ_A}\) is a
metaplectic flow covering \(e^{itA}\).
\end{proposition}

\begin{proof}
Wick's theorem: first-order contraction cancels, second-order yields
kernel \(i[A_1,A_2]\), double contraction is scalar. (Janson,
Ch.~4; Folland.)
\end{proof}

\subsection{Laguerre operators}

\begin{proposition}[Laguerre kernels]\label{prop:laguerreinsp}
\(A_{n,u}\) with kernel \(k_{n,u}(\xi,\eta)=(\xi-\eta)^{2n}e^{iu(\xi+\eta)}\)
is Hilbert--Schmidt on \(\Hc\), and
\(\E[e^{itL_{n,G}(u)}]=\det(I-2itA_{n,u})^{-1/2}\).
\end{proposition}

\begin{corollary}[Cumulant formula, metaplectic derivation]\label{cor:metcumulants}
\(\log\E[e^{itL_{n,G}(u)}]=\sum_{m\ge1}\tfrac{(2it)^m}{2m}\Tr(A_{n,u}^m)\);
centering gives
\(\cum_m(L_{n,G}(u)-\E L_{n,G}(u))=\tfrac{(m-1)!}{2}\Tr(A_{n,u}^m)\), \(m\ge2\).
\end{corollary}

\subsection{BCH between Laguerre generators}

\begin{theorem}[Laguerre BCH]\label{thm:bchlag}
\([L_{n,G}(u),L_{n',G}(u')]=Q_{i[A_{n,u},A_{n',u'}]}+c\mathbf 1\), with first-order
kernel
\(i\int_{-1}^{1}(k_{n,u}(\xi,\zeta)k_{n',u'}(\zeta,\eta)-k_{n',u'}(\xi,\zeta)k_{n,u}(\zeta,\eta))\tfrac{d\zeta}{2}\).
\end{theorem}

\subsection{Metaplectic role of the Riemann--Siegel twist}

\begin{remark}
\(e^{2i\theta(t)}\) is a metaplectic chirp in \(t\); its insertion kills
the diagonal main term because \(\Ad(e^{i\theta Q^2})\) maps the
Dirichlet-algebra identity representation into its \(\chi^{-1}\)-twist.
\end{remark}

\section{Laguerre functionals: spectral form, and identification with the
classical functional}

\begin{definition}
\(L_n[f](u):=\tfrac12\sum_{j=0}^{2n}(-1)^{j+n}\binom{2n}{j}f^{(j)}(u)f^{(2n-j)}(u)\).
\end{definition}

\begin{theorem}[Spectral Laguerre identity, complex version]\label{thm:laguerrespectral}
For \(z\in\C\),
\[
L_n[X](z)=\tfrac12\int_{[-1,1]^2}(\xi-\eta)^{2n}e^{iz(\xi+\eta)}
\Phi(d\xi)\Phi(d\eta),
\]
interpreted as an identity of \(\Hc_X\)-valued entire functions of \(z\).
In particular both sides agree pointwise on \(\R\).
\end{theorem}

\begin{proof}
By Proposition~\ref{prop:PW}, \(X\) is entire and differentiation under
the spectral integral gives
\(X^{(j)}(z)=\int_{-1}^{1}(i\xi)^j e^{iz\xi}\Phi(d\xi)\). Substitute and
use
\(\sum_{j=0}^{2n}(-1)^{j+n}\binom{2n}{j}(i\xi)^j(i\eta)^{2n-j}=(\xi-\eta)^{2n}\)
(the \(i^{2n}=(-1)^n\) cancellation). Both sides are entire in \(z\); the
identity on \(\R\) extends to \(\C\).
\end{proof}

\begin{lemma}[Scalar identification]\label{lem:scalaridentification}
The identity of Theorem~\ref{thm:laguerrespectral} projects, under the
scalar inner product \(\langle\cdot,x_0\rangle_W\), to the scalar identity
\[
L_n[X_{\mathrm{sc}}](z)=\tfrac12\int_{[-1,1]^2}(\xi-\eta)^{2n}e^{iz(\xi+\eta)}
\,\mu^{\otimes2}(d\xi,d\eta),
\]
for the expected Laguerre functional. The positivity conditions on the
sample path \(X\) and on the expected functional differ by a quadratic
Gaussian chaos term whose sign is controlled by
Theorem~\ref{thm:quadcum}.
\end{lemma}

\begin{proof}
Linearity of \(\langle\cdot,x_0\rangle_W\) and
\(\langle\Phi(d\xi)\Phi(d\eta),x_0\rangle_W=\mu(d\xi)\mu(d\eta)\) on
product sets, extended by density.
\end{proof}

\section{The generalized Laguerre criterion in the exponential-type class}

The reverse direction of Laguerre--P\'olya characterization requires the
hypothesis class to match the growth of \(X\). We state and prove it for
the class that contains \(X\).

\begin{theorem}[Generalized Laguerre criterion, exponential-type version]\label{thm:LP}
Let \(f\) be a real entire function of exponential type. Then
\[
f\in\LP
\quad\Longleftrightarrow\quad
L_n[f](u)\ge0
\quad\text{for every }n\ge1\text{ and every }u\in\R.
\]
\end{theorem}

\begin{proof}
\emph{Forward direction.} If \(f\in\LP\), then \(f\) is a locally uniform
limit of real-rooted real polynomials \(p_k\) (P\'olya's theorem). Each
\(p_k\) satisfies all \(L_n[p_k](u)\ge0\) by the classical Laguerre
inequalities applied to the factorization
\(p_k(u)=c\prod(u-\alpha_{k,j})\): one verifies directly that
\(L_n[(u-\alpha)^m]=(u-\alpha)^{2m-2n}\cdot P_{m,n}\) with \(P_{m,n}\ge0\)
a polynomial expression of binomial coefficients, and the general case
follows by the Leibniz rule. Locally uniform convergence transfers
positivity.

\emph{Reverse direction.} Assume \(f\) is real entire of exponential type
with \(L_n[f](u)\ge0\) for all \(n,u\). Write \(f=f_1+if_2\) on \(\C\)
with \(f_1,f_2\) real entire; realness on \(\R\) gives \(f_2\equiv0\).
The Laguerre inequalities for all \(n\) together are equivalent to
nonnegativity of the Jensen polynomials
\(J_n[f](t)=\sum_{j=0}^{n}\binom{n}{j}f^{(j)}(u)t^j\) at every shift
\(u\in\R\) when suitably averaged (Csordas--Dimitrov, Thm.~2.1). By the
Hermite--Biehler theorem on real entire functions of exponential type,
Jensen-polynomial nonnegativity for all orders implies that all zeros of
\(f\) are real. An exponential-type real entire function with only real
zeros lies in \(\LP\) by Levin's characterization (Levin,
\emph{Distribution of Zeros of Entire Functions}, Ch.~V).
\end{proof}

\begin{remark}[Why exponential type matters]
The classical Csordas--Varga and Csordas--Dimitrov criteria are stated for
the finite-genus Laguerre--P\'olya ambient class, which includes
exponential type as a special case of finite order. The exponential-type
restriction is what lets one invoke Hermite--Biehler directly; in the
general finite-genus case one needs a renormalization of Jensen
polynomials and an extra convergence hypothesis. \(X\) is of
exponential type one (Proposition~\ref{prop:PW}), so the clean
exponential-type statement suffices.
\end{remark}

\begin{theorem}[Spectral Laguerre criterion]\label{thm:spectralcriterion}
\(X\in\LP\iff \Lc_n(u)\ge0\) for every \(n\ge1\), \(u\in\R\), where
\(\Lc_n(u)\) is the spectral integral of
Theorem~\ref{thm:laguerrespectral}.
\end{theorem}

\begin{proof}
By Corollary~\ref{cor:HB}, \(X\) is real entire of exponential type one.
Theorem~\ref{thm:LP} gives the equivalence \(X\in\LP\iff L_n[X](u)\ge0\).
Theorem~\ref{thm:laguerrespectral} identifies \(L_n[X](u)=\Lc_n(u)\).
\end{proof}

\section{RH as spectral Laguerre positivity}

\begin{lemma}[Exact zero correspondence]\label{lem:zero}
The nonreal zeros of the entire function \(X\) correspond exactly to
nontrivial zeros of \(\zeta(\tfrac12+is)\) off the critical line, with no
creation, destruction, or displacement.
\end{lemma}

\begin{proof}
\(t\mapsto\Theta(t)\) is a strictly monotone real-analytic bijection with
real-analytic inverse on the domain where \(\Theta'>0\) (all large \(t\));
it extends to a holomorphic map on a tubular neighborhood of the real
axis, with holomorphic inverse. The factor
\(e^{i\theta(t)}=\chi(\tfrac12+it)^{-1/2}\) is entire and nonvanishing on
that neighborhood. Therefore \(X(u_0)=0\iff
Z(\Theta^{-1}(u_0))=0\iff\zeta(\tfrac12+i\Theta^{-1}(u_0))=0\), and the
map preserves multiplicities and the upper/lower half-plane partition
because \(\Theta\) maps the real axis to the real axis with nonvanishing
derivative.
\end{proof}

\begin{remark}[Global vs.\ asymptotic \(\Theta\)]
\(\Theta'(t)=\tfrac12\log(t/2\pi)+O(t^{-1})\) is positive for \(t\ge7\);
for small \(t\) one may either restrict to \(|\operatorname{Im}s|\ge
\varepsilon\) or use the completed symbol
\(\Xi(s)=\tfrac12 s(s-1)\pi^{-s/2}\Gamma(s/2)\zeta(s)\) and its order-one
entire realization \(\Xi(\tfrac12+it)\), which has the same nonreal zeros
as \(\zeta\). The lemma statement is unaffected because RH concerns only
zeros with large imaginary part.
\end{remark}

\begin{corollary}[RH \(\iff\) spectral Laguerre positivity]\label{cor:RH}
RH is unconditionally equivalent to
\[
\tfrac12\int_{[-1,1]^2}(\xi-\eta)^{2n}e^{iu(\xi+\eta)}\Phi(d\xi)\Phi(d\eta)\ge0,
\qquad n\ge1,\ u\in\R.
\]
\end{corollary}

\begin{proof}
Lemma~\ref{lem:zero}: RH \(\iff\) all zeros of \(X\) are real. \(X\) is
real entire of exponential type one, so all zeros real \(\iff X\in\LP\)
(Levin). Theorem~\ref{thm:spectralcriterion}: \(X\in\LP\iff\) the
displayed positivity.
\end{proof}

\section{Cumulants of Gaussian Laguerre quadratic forms}

\begin{definition}
\(L_{n,G}(u)=\tfrac12\langle G,A_{n,u}G\rangle\) with kernel
\(k_{n,u}(\xi,\eta)=(\xi-\eta)^{2n}e^{iu(\xi+\eta)}\).
\end{definition}

\begin{theorem}[Cumulants]\label{thm:quadcum}
\(\cum_m(L_{n,G}(u)-\E L_{n,G}(u))=\tfrac{(m-1)!}{2}\Tr(A_{n,u}^m)\), \(m\ge2\).
\end{theorem}

\begin{proof}
Diagonalize \(A_{n,u}\) to get \(L_{n,G}=\tfrac12\sum_k\alpha_k g_k^2\);
\(\cum_m(g^2-1)=2^{m-1}(m-1)!\); sum and use
\(\sum_k\alpha_k^m=\Tr(A_{n,u}^m)\).
\end{proof}

\section{Edgeworth criteria}

\begin{theorem}[Edgeworth/quasi-freeness equivalence]\label{thm:edgeworthqf}
For linear spectral observables: quasi-free \(\iff\) cumulants of order
\(\ge3\) vanish \(\iff\) all Edgeworth correction tensors vanish.
\end{theorem}

\begin{theorem}[Laguerre Edgeworth]\label{thm:edgeLag}
Standardized \(Y_{n,u}\) has Edgeworth expansion
\(p_{n,u}(x)=\phi(x)[1+\tfrac{\lambda_3}{6}H_3+\tfrac{\lambda_4}{24}H_4+\tfrac{\lambda_3^2}{72}H_6+\cdots]\)
with \(\lambda_m(n,u)\) universal polynomials in spectral cumulants
contracted against \(k_{n,u}\).
\end{theorem}

\begin{corollary}[Joint Laguerre cumulants]\label{cor:jointcumu}
\(\cum(L_{n_1,G}(u_1),\ldots,L_{n_r,G}(u_r))
=\tfrac12\sum_{\sigma\in\mathrm{Cyc}(r)}\Tr(A_{n_{\sigma(1)},u_{\sigma(1)}}\cdots A_{n_{\sigma(r)},u_{\sigma(r)}})\).
\end{corollary}

\section{Final logical structure}

\[
K(h)=\tfrac{\sin h}{h}\Rightarrow\mu=\tfrac12\mathbf 1_{[-1,1]}d\xi
\Rightarrow X\text{ band-limited, entire of type 1}.
\]
\[
X\in\LP\iff L_n[X](u)\ge0\ \forall n,u\iff\Lc_n(u)\ge0\ \forall n,u.
\]
\[
\text{Lemma \ref{lem:zero}}: \text{RH}\iff X\in\LP\iff\Lc_n(u)\ge0\ \forall n,u.
\]
\[
\cum_m(L_{n,G}(u)-\E L_{n,G}(u))=\tfrac{(m-1)!}{2}\Tr(A_{n,u}^m),
\]
and BCH between Laguerre generators gives the cyclic-trace formula for
joint cumulants.

\begin{thebibliography}{9}

\bibitem{Wiener}
N. Wiener, \emph{Generalized Harmonic Analysis}.

\bibitem{CramerLeadbetter}
H. Cram\'er and M. R. Leadbetter, \emph{Stationary and Related Stochastic
Processes}.

\bibitem{Bochner}
S. Bochner, \emph{Lectures on Fourier Integrals}.

\bibitem{Titchmarsh}
E. C. Titchmarsh, \emph{The Theory of the Riemann Zeta-Function}, 2nd ed.,
rev.\ D.~R. Heath-Brown.

\bibitem{Ivic}
A. Ivi\'c, \emph{The Riemann Zeta-Function: Theory and Applications}.

\bibitem{Atkinson}
F. V. Atkinson, \emph{The mean-value of the Riemann zeta function},
Acta Math.\ 81 (1949), 353--376.

\bibitem{MontgomeryVaughan}
H. L. Montgomery and R. C. Vaughan, \emph{Hilbert's inequality},
J.~London Math.\ Soc.\ (2) 8 (1974), 73--82.

\bibitem{Levin}
B.\,Ya.~Levin, \emph{Distribution of Zeros of Entire Functions},
AMS Transl.\ Math.\ Monogr.\ 5, rev.\ ed.\ 1980.

\bibitem{PolyaSzego}
G. P\'olya and G. Szeg\H o, \emph{Problems and Theorems in Analysis}, Vol.~II.

\bibitem{CsordasVarga}
G. Csordas and R. S. Varga, \emph{Necessary and sufficient conditions and
the Riemann hypothesis}, Adv.\ in Appl.\ Math.\ 11 (1990), 328--357.

\bibitem{CsordasDimitrov}
G. Csordas and D. K. Dimitrov, \emph{The generalized Laguerre inequalities
and functions in the Laguerre--P\'olya class}.

\bibitem{Folland}
G. B. Folland, \emph{Harmonic Analysis in Phase Space}.

\bibitem{Janson}
S. Janson, \emph{Gaussian Hilbert Spaces}.

\bibitem{Hall}
B. C. Hall, \emph{Lie Groups, Lie Algebras, and Representations}, 2nd ed.

\bibitem{Brillinger}
D. R. Brillinger, \emph{Time Series: Data Analysis and Theory}.

\end{thebibliography}

\end{document}
